
\section{Input variables} \label{ch:inputvar}

\textcolor{red}{
{\it The old input file has now been transformed to fortran NAMELIST format, which allows key=value pair reading like for an .ini file; the next subsection describes the new NAMELIST input format; while most input data stay the same (only the form of the input data file changes), some new data are appearing, which can now be changed directly using input. Previously these data had been compiled into the code. Also note that now initial and restart runs have the {\bf same} input files.}}

The input control file of NBODY6++ (see below), contains
order 90 parameters which guide one simulation run for
its technical and physical properties
(it is very similar but not identical to the one used for NBODY6).
As for the technical aspect, the file supervises the run e.g.
for its duration, intervals of the output, or error check;
the physical parameters concern the size of a cluster,
initial conditions, or a number of optional features 
related to the numerical problem to be studied.
The handling of this input file appears rather entangled
at first sight, for it has grown rather historically and
``ready--for--use'' than custom--oriented.
Thus, the input variables are read by different routines
(functions) in the code, and the nature of the parameters
are woven with each other in some cases.
Also, some parameters require additional input, such that the
total number of lines and parameters may vary.

In the following, we explain the main input file and give an
example of typical values for a simulation of an isolated
globular cluster.
Then, we proceed to the thresholds.\\

\noindent

{\tt\scriptsize
\begin{longtable}{|p{0.6in}|*{10}{l}|} 
\caption{\texttt{Input with all options}}
\label{table:allinput}\\
\hline
nbody6.F & \textcolor{red}{Input}& \textcolor{red}{Namelist}&  \textcolor{red}{INNBODY6} &&&&&&&\\
         & KSTART   & TCOMP    & TCRTP0   & isernb   & iserreg  
         & iserks   &          &          &          &          \\\hline
input.F &  \textcolor{red}{Input}& \textcolor{red}{Namelist}&  \textcolor{red}{ININPUT} &&&&&&&\\
        & N        & NFIX     &  NCRIT   & NRAND    & NNBOPT   &
           NRUN    & NCOMM    &          &          &          \\

        & ETAI     & ETAR     & RS0      & DTADJ    & DELTAT   &
           TCRIT    & QE       & RBAR     & ZMBAR    &          \\

        & KZ(1)    & KZ(2)    & KZ(3)    & KZ(4)    & KZ(5)    &
           KZ(6)    & KZ(7)    & KZ(8)    & KZ(9)~~~~KZ(10)&   \\
         
        & KZ(11)   & KZ(12)   & KZ(13)   & KZ(14)   & KZ(15)   &
           KZ(16)   & KZ(17)   & KZ(18)   & KZ(19)~~~KZ(20)&   \\
         
        & KZ(21)   & KZ(22)   & KZ(23)   & KZ(24)   & KZ(25)   &
           KZ(26)   & KZ(27)   & KZ(28)   & KZ(29)~~~KZ(30)&   \\
         
        & KZ(31)   & KZ(32)   & KZ(33)   & KZ(34)   & KZ(35)   &
           KZ(36)   & KZ(37)   & KZ(38)   & KZ(39)~~~KZ(40)&   \\
         
        & KZ(41)   & KZ(42)    & KZ(43)    & KZ(44)    & KZ(45)    &
          KZ(46)    & KZ(47)    & KZ(48)    & KZ(49)~~~KZ(50)&   \\
         
        & DTMIN    & RMIN     & ETAU     & ECLOSE   & GMIN     &
           GMAX     & SMAX     &          &          &          \\\hline

data.F & \textcolor{red}{Input}& \textcolor{red}{Namelist}&  \textcolor{red}{INDATA} &&&&&&&\\
       & ALPHA    & BODY1    & BODYN    & NBIN0    & NHI0     &
          ZMET     & EPOCH0   & DTPLOT   &          &          \\\hline
setup.F & \textcolor{red}{Input}& \textcolor{red}{Namelist}&  \textcolor{red}{INSETUP} &&&&&&&\\
         & AP0      & ECC      & N2       & SCALE   &       
         &          &          &          &(KZ(5)$=$2)&\\
         & AP0      & ECC      & SCALE    &         & 
         &          &          &          &(KZ(5)$=$3)&\\
         & AP0      & ECC      & SCALE    &         & 
&          &          &          &(KZ(5)$=$3)&\\
         & SEMI     & ECC      & M1       & M2      &
         &          &          &          &    (KZ(5)$=$4)&\\
         & ZMH      & RCUT     &          &         &
          \multicolumn{5}{r}{(KZ(5)$=$6\&\&KZ(24)<0)}&\\\hline

scale.F & \textcolor{red}{Input}& \textcolor{red}{Namelist}&  \textcolor{red}{INSCALE} &&&&&&&\\
         & Q        & VXROT    & VZROT    & RTIDE   &          
         &          &          &          &         &          \\ \hline

xtrnl0.F & \textcolor{red}{Input}& \textcolor{red}{Namelist}&  \textcolor{red}{INXTRNL0} &&&&&&&\\
         & GMG      & RG0      &          &         &
         &          &          &          &       (KZ(14)$=$2) & \\
         & GMG      & DISK~~~A &
         \multicolumn{6}{l}{B~~~VCIRC~~~RCIRC~~~GMB~~~
          AR~~~GAM~~~RG[1:3]~~~VG[1:3]} &       (KZ(14)$=$3) & \\
         & MP       & AP       & MPDOT    & TDELAY  &
         &           \multicolumn{4}{r}{(KZ(14)$=$3||KZ(14)$=$4)} & \\ \hline

binpop.F & \textcolor{red}{Input}& \textcolor{red}{Namelist}&  \textcolor{red}{INBINPOP} &&&&&&&\\
         & SEMI     & ECC      & RATIO    & RANGE   & NSKIP
         & IDORM    & \multicolumn{3}{r}{(KZ(8)$=$1||KZ(8)$>$4)} & \\ \hline

hipop.F & \textcolor{red}{Input}& \textcolor{red}{Namelist}&  \textcolor{red}{INHIPOP} &&&&&&&\\
         & SEMI     & ECC      & RATIO    & RANGE   & 
         &          &\multicolumn{3}{r}{(KZ(8)$>$0$\&\&$KZ(18)$>$1)}&\\ \hline
imbhinit.F & \textcolor{red}{Currently} & \textcolor{red}{not} & \textcolor{red}{supported} & \textcolor{red}{by} & \textcolor{red}{NAMELIST} & \textcolor{red}{Input} &&&&\\
         & MMBH  & XBH(1)   & XBH(2)    & XBH(3)  & VBH(1)
         & VBH(2)   & VBH(3)   & DTBH     &    (KZ(24)$=$1) & \\\hline

cloud0.F & \textcolor{red}{Currently} & \textcolor{red}{not} & \textcolor{red}{supported} & \textcolor{red}{by} & \textcolor{red}{NAMELIST} & \textcolor{red}{Input} &&&& \\
cloud0.F & NCL      & RB2      & VCL      & SIGMA   & CLM 
         & RCL2     &          &          &    (KZ(13)$>$0) & \\\hline
\end{longtable}
}


\vspace{0.2cm}
\noindent
\begin{longtable}{@{}p{1.5cm}p{13.0cm}}
\caption{\texttt{Input of \&INNBODY6, nbody6.F}}
\label{table:innbody6}\\\hline
%
KSTART  & Run control index\\
        & =1: new run (construct new model or read from \texttt{dat.10}) \\
        & =2: restart/continuation of a run, needs \texttt{comm.1}\footnote{the user should create a comm.1 file by copy-pasting the comm.2\_X, where X is the wanted NBU time value to restart the run. Notice that the same applies to KSTART $>$ 2; notice also that \\
        \texttt{cp -p `ls -rt comm.2* | tail -1` comm.1} \\
        is a nice unix way to automatically copy the last comm.2 to comm.1 }\\
%
TCOMP   & Maximum wall-clock time in seconds (parallel runs: wall clock)\\
%
TCRTP0  & Termination time in Myr \\
%
isernb  & For MPI parallel runs: only irregular block sizes larger than
          this value are executed in parallel mode
          (dummy variable for single CPU) \\
iserreg & For MPI parallel runs: only regular block sizes larger than
          this value are executed in parallel mode
          (dummy variable for single CPU) \\
iserks  & For MPI parallel runs: only ks block sizes larger than this value
           are executed in parallel mode 
          (dummy variable for single CPU) \\

\end{longtable}

\hrule
\noindent
\begin{longtable}{@{}p{1.5cm}p{13.0cm}}
\caption{\texttt{Input of \&ININPUT, input.F line 1}}
\label{table:ininput1}\\\hline
%
N       & Total number of particles (single + c.m.s.~of binaries;
          singles + 3$\times$c.m.s.~of binaries < NMAX$-$2)\\
NFIX    & Multiplicator for output interval of data on \texttt{conf.3} and
          of data for binary stars (output each DELTAT$\times$NFIX time steps; compare KZ(3) and KZ(6))\\
NCRIT   & Minimum particle number (alternative termination criterion) \\
NRAND   & Random number seed; any positive integer \\
%          (< LMAX$-$5) set in \texttt{params.h}\\
NNBOPT  & Desired optimal neighbour number (< LMAX$-$5)\\ %(R.Sp.)
NRUN    & Run identification index\\
NCOMM   & Multiplicator for frequency to store the dumping data (mydump); positive integer; basic frequency determined by DTADJ \\
\end{longtable}

\hrule
\noindent
\begin{longtable}{@{}p{1.5cm}p{13.0cm}}
\caption{\texttt{Input of \&ININPUT, input.F line 2}}
\label{table:ininput2}\\\hline
%
ETAI    & Time--step factor for irregular force polynomial\\
ETAR    & Time--step factor for regular force polynomial\\
RS0     & Initial guess for all radii of neighbour spheres ($N$--body units)\\
DTADJ   & Time interval for parameter adjustment and energy check ($N$--body units). In a good practice, DTADJ should be the same as, or larger than, DELTAT. \\
%%          (in TCR if KZ(35)=0, otherwise in $N$-body units)\\
DELTAT  & Time interval for writing output data and diagnostics,
          multiplied by NFIX ($N$--body units).\\
%%        & -- in units of $t_{\rm cr}$ if KZ(35) = 0;\\
%%$        & -- in scaled units if KZ(35) > 0.\\
%%        & for DTADJ=DELTAT (recommended) output data is written every
%%          adjust time\\
TCRIT   & Termination time ($N$--body units) \\
%*->             The _earlier_ termination criterion becomes active
QE      & Energy tolerance:\\
        & -- immediate termination if DE/E > 5*QE \& KZ(2) $\leq$ 1;\\
        & -- restart if DE/E > 5*QE \& KZ(2) > 1 and
         termination after second restart attempt.\\
RBAR    & Scaling unit in pc for distance ($N$--body units)\\
ZMBAR   & Scaling unit for average particle mass in solar masses\\
        & (in scale-free simulations RBAR and ZMBAR can be set to zero; depends on KZ(20)) \\
\end{longtable}

\hrule
\noindent
\begin{longtable}{@{}p{1.5cm}p{13.0cm}}
\caption{\texttt{Input of \&ININPUT, input.F lines 3-7}}
\label{table:ininput37}\\\hline
%
KZ(1)   & Save COMMON to file \texttt{comm.1}\\
        & $=1$: at end of run or when dummy file STOP is created\\
        & $=2$: every 100*NMAX steps\\
KZ(2)   & Save COMMON to file \texttt{comm.2}\\
        & $=1$: save at output time\\
        & $=2$: save at output time and restart simulation if energy error DE/E > 5*QE\\
KZ(3)   & Save basic data to file \texttt{conf.3} at output time (unformatted)\\
KZ(4)   & (Suppressed) Binary diagnostics on \texttt{bdat.4} (\# = threshold levels <10)\\
KZ(5)   & Initial conditions of the particle distribution, needs KZ(22)=0\\
        & $=0$: uniform \& isotropic sphere\\
        & $=1$: Plummer random generation\\
        & $=2$: two Plummer models in orbit (extra input)\\
        & $=3$: massive perturber and planetesimal disk (each pariticle 
              has circular orbit, constant separation along 
              radial direction between each neighbor and random phase)
              (extra input)\\
        & $=4$: massive initial binary (extra input)\\
        & $=5$: Jaffe model (extra input) \\
        & $\ge 6$: Zhao BH cusp model (extra input if KZ(24)<0) \\
KZ(6)   & Output of significant and regularized binaries at main output (\texttt{bodies.f})\\
        & $=1$: output regularized and significant binaries (|E|>0.1 ECLOSE)\\
        & $=2$: output regularized binaries only\\
        & $=3$: output significant binaries at output time and regularized binaries with time interval DELTAT\\
        & $=4$: output of regularized binaries only at output time\\
KZ(7)   & Determine Lagrangian radii and average mass, particle counters, average velocity, velocity dispersion, rotational velocity within Lagrangian radii (\texttt{lagr.f})\\
        & $=1$: Get actual value of half mass radius RSCALE by using current total mass\\
        & $\ge 2$: Output data at main output and \texttt{lagr.7} \\
        & $\ge 6$: Output Lagrangian radii for two mass groups at \texttt{lagr.31} and \texttt{lagr.32} (\texttt{lagr2.f}; based on KZ(5)=1,2; cost is O($N^2$))\\
        & ---- methods:\\
        & $=2,4$: Lagrangian radii calculated by initial total mass \\
        & $=3,\ge 5$: Lagrangian radii calculated by current total mass (The single/K.S-binary Lagrangian radii are still calculated by initial single/binary total mass)\\
        & $=2,3$: All parameters are averaged within the shell between two Lagrangian radii neighbors \\
        & $\ge 4$: All parameters are averaged from center to each Lagrangian radius\\
KZ(8)   & Primordial binaries initialization and output (\texttt{binpop.f})\\
        & ---- Initialization:\\
        & $=0$: No primordial binaries \\
        & $=1,\ge 3$: generate primordial binaries based on KZ(41) and KZ(42) (\texttt{binpop.F}) \\
        & $=2$: Input primordial binaries from first 2$\times$NBIN0 lines of \texttt{dat.10} \\
        & ---- Output:\\
        & $>0$: Save information of primordial binary that change member in \texttt{pbin.18}; binary diagnostics at main output (\texttt{binout.f}) \\
        & $\ge 2$: Output KS binary in \texttt{bdat.9}, soft binary in \texttt{bwdat.19} at output time\\
KZ(9)   & Binary diagnostics \\
        & $=1,3$: Output diagnostics for the hardest binary below ECLOSE in \texttt{hbin.39} (\texttt{adjust.f}) \\
        & $\ge 2$: Output binary evolution stages in \texttt{binev.17} (\texttt{binev.f}) \\
        & $\ge 3$: Output binary with degenerate stars in \texttt{degen.4} (\texttt{degen.f}) \\
KZ(10)  & K.S. regularization diagnostics at main output \\
        & $>0$: Output new K.S. information \\
        & $>1$: Output end K.S. information \\
        & $\ge 3$: Output each integrating step information\\
KZ(11)  & (Suppressed) \\
% KZ(11)  & Hierarchical systems (\texttt{hiarch.f})\\
%        & $=1,3$: write to \texttt{hia.12} \\
%        & $=2$: create primordial triples (\texttt{hipop.f})\\
%        & $=3$: both \\
KZ(12)  & $>0$: HR diagnostics of evolving stars with output time interval DTPLOT in \texttt{sse.83} (single star) and \texttt{bse.82} (K.S. binary) \\
        & $=-1$: used if KZ(19)$=0$ (see details in KZ(19) description)\\
KZ(13)  & Interstellar clouds \\
        & $=1$: constant velocity for new cloud\\
        & $>2$: Gaussian velocity for new cloud\\
KZ(14)  & External tidal force \\
        & $=0$: no external tidal field\\
        & $=1$: standard solar neighbor tidal field\\
        & $=2$: point-mass galaxy with circular orbit (extra input) \\
        & $=3$: point-mass + disk + halo + Plummer (extra input) \\
        & $=4$: Plummer model (extra input) \\
        & $=5$: Milky-Way potential (extra input) based on Galpy package (Bovy, 2015) \\
        & Note: Namelist input has all fields for values 2,3,4, but user must take care to put proper values \\
KZ(15)  & Triple, quad, chain and merger search. Note that \textit{triple} and \textit{quad} in KZ(15) means unperturbed triple and quad systems. \\
        & $\ge 1$: Switch on triple, quad, chain (KZ(30)>0) and merger search (\texttt{impact.f}) \\
        & $\ge 2$: Diagnostics at main output at begin and end of triple, quad \\
        & $\ge 3$: Save first five outer orbits every half period of wide quadruple before merger and stable quadruples accepted for merger in \texttt{quastab.89} \\
KZ(16)  & Auto-adjustment of regularization parameters. WARNING: not tested for large $N$ ($N>10^5$) systems with large mass range. \\
        & $\ge 1$: Adjust RMIN, DTMIN \& ECLOSE every DTADJ time \\
        & $\ge 3$: modify RMIN for GPERT > $0.05$ or < $0.002$ in chain; output diagnostics at \texttt{kscrit.77} \\
KZ(17)  & Auto-adjustment of ETAI, ETAR and ETAU by tolerance QE every DTADJ time (\texttt{check.f}). WARNING: not tested for large $N$ ($N>10^5$) systems with large mass range. \\
        & $\ge 1$: Adjust ETAI, ETAR \\
        & $\ge 2$: Adjust ETAU \\
KZ(18)  & Hierarchical systems \\
        & $=1,3$: diagnostics (\texttt{hiarch.f})\\
        & $\ge 2$: Initialize primordial stable triples, number is NHI0 (\texttt{hipop.F}) \\
        & $\ge 4$: Data bank of \textit{hierarchical} stable triple, quad in \texttt{hidat.87} (\texttt{hidat.f}) \\
KZ(19)  & Stellar evolution and mass loss scheme \\
        & $=0$: if KZ(12)$=-1$, the output data will keep the input data unit if KZ(22)$=2-4$ or $N$-body units if KZ(22)$=6-10$\\
        & $=1,2$: supernova scheme \\
        & $\ge 3$: Eggleton, Tout \& Hurley \\
        & $\ge 5$: extra diagnostics ({mdot.F}) \\
        & $=2,4$: input stellar parameters from \texttt{datsev.21}; 
                  read by (\texttt{instar.f}): \\
                & N lines of (MI, KW, M0, EPOCH1, OSPIN, MC); N is the number of particles in \texttt{datsev.21} and \texttt{dat.10} and has to be set accordingly in Nbody input file. \\
                & ~~~~~~   MI: current mass (solar mass) \\
                & ~~~~~~   KW: stellar type \\
                & ~~~~~~   M0: initial mass (solar mass) \\
                & ~~~~~~   EPOCH1: evolved age of star (AGE $= -$ EPOCH1 [Myr]) \\
                & ~~~~~~   OSPIN: angular velocity of star (may be zero if not relevant) \\
                & ~~~~~~   MC: core mass (solar mass, may be zero if not relevant) \\
        & $=4$: Take input data from other code (e.g. Monte Carlo); \\ 
        & first line of \texttt{datsev.21} in this case read by \texttt{start.F} to establish desired scaling: \\
                & 1 line of (TMOC,NZERO,RBAR,ZMBAR,TSTAR) \\ 
                &  \ff{(these data override data read or computed by Nbody before)} \\
                & ~~~~~~   TMOC: Cluster age in Myr \\
                & ~~~~~~   NZERO: Particle number at zero cluster age \\
                & ~~~~~~   RBAR, ZMBAR: scaling for radii and average \\  & ~~~~~~   TSTAR: 1Myr in model units; \\
KZ(20)  & Initial mass functions, need KZ(22)$=$0 or 9: \\
        & $=0$: self-defined power-law mass function using ALPHAS (\texttt{data.F}) \\
        & $=1$: Miller-Scalo-(1979) IMF (\texttt{imf.f}) \\
        & $=2,4$: KTG (1993) IMF (\texttt{imf2.f}) \\
        & $=3,5$: Eggleton-IMF (\texttt{imf2.f}) \\
        & $=6,7$: Kroupa(2001) (\texttt{imf2.f}), extended to Brown Dwarf regime (\texttt{imfbd.f}) \\
        & ---- Primordial binary mass \\
        & $=2,6$: random pairing (\texttt{imf2.f}) \\
        & $=3,4,5,7$: binary mass ratio corrected by $(m_1/m_2)\prime = (m_1/m_2)^{0.4}$ + constant (Eggleton, \texttt{imf2.f}) \\
        & $=8$: binary mass ratio $q=m_1/m_2$ ($m_2 \le m_1$) use distribtution $0.6 q^{-0.4}$ (Kouwenhoven) \\
%        & $=7$: Distinct mass bins  (\texttt{imf3.f}) \\
%        & $=8$: Scalo (1986)-IMF  with new random function RAN2 (\texttt{imf3.f}) \\
KZ(21)  & Extra diagnostics information at main output every DELTAT interval (\texttt{output.F}) \\
        & $\ge 1$: output NRUN, MODEL, TCOMP, TRC, DMIN, AMIN, RMAX, RSMIN, NEFF \\
        & $\ge 2$: Number of escapers NESC at main output will be counted by Jacobi escape criterion (cost is O($N^2$), \texttt{jacobi.f}) \\
KZ(22)  & Initialization of basic particle data mass, position and velocity (\texttt{data.F}) \\
        & ---- Initialization with internal method \\
        & $=0,1$: Initial position, velocity based on KZ(5), initial mass based on KZ(20) \\
        & $=1$: write initial conditions in \texttt{dat.10} (\texttt{scale.F}) \\
        & ---- Initialization by reading data from \texttt{dat.10} \\
        & $=2$: input through NBODY-format (7 parameters each line: mass, position(1:3), velocity(1:3))\\
        & $=3$: input through Tree-format (\texttt{data.F}) \\
        & $=4$: input through Starlab-format \\
        & $=6$: input through NBODY-format and do scaling\\
        & $=7$: input through Tree-Format and do scaling\\
        & $=8$: input through Starlab-format and do Scaling\\
        & $=9$: input through NBODY-format but ignore mass (first column) and use IMF based on KZ(20), then do scaling \\
        & $=10$: input through NBODY-format and all units are astrophysical units (mass: $M_\odot$; position: pc; velocity: km/s)\\
KZ(23)  & Removal of escapers (\texttt{escape.F})\\
        & $= 0, \ge 5$: no removal of escapers \\ 
        & $\ge 1$: remove escapers, and ghost particles generated by stellar evolution or two star coalescence (collision) \\
        & $=2,3,4$: write escaper diagnostics in \texttt{esc.11} and \texttt{escbin.31} for binary systems. \\
        & $=3$: tidal cutoff escape criterion (this option has been misused for tidal tail for a while) \\
        & $=4$: escape angles in main output \\
        & $\ge 5$: initialization \& integration of tidal tail (was wrongly: $\ge 3$)\\
%        & $=4$: write \#3 to diagnostics \\
KZ(24)  & =1: should not be used anymore, not supported (was: subsystems and IMBH)\\
        & =2: start with only black holes at T=0. \\
%KZ(25)  & Partial reflection of KS binary orbit (GAMMA < GMIN; suppressed) \\
KZ(25)  & Velocity kicks for white dwarfs (\texttt{kick.F}) \\
        & $=1$: Type 10 Helium white dwarf \& 11 Carbon-Oxygen white dwarf \\
        & $=2$: All WDs (type 10, 11 and type 12 Oxygen-Neon white dwarf) \\
KZ(26)  & Slow-down of two-body motion, increase the regularization integration efficiency \\
        & $\ge 1$: Apply to KS binary \\
        & $\ge 2$: Apply to chain \\
        & $=3$: Rectify to get better energy conservation \\
KZ(27)  & Two-body tidal circularization (Mardling \& Aarseth, 2001; Portegies Zwart et al. 1997) (for values 1,2) plus orbit shrinking of relativistic compact binaries (Rizzuto et al. 2021, 2022; Arca Sedda et al. 2021 (for values 3,4).
        (Please suppress in KS parallel version) \\
        & $=1$: sequential  \\
        & $=2$: chaos \\
        & $=3$: GR energy loss with sequential tides\\
        & $=4$: GR energy loss with chaos tides \\
        & $=-1$: Only detect collision and suppress coalescence \\
KZ(28)  & Magnetic braking. Gravitational radiation for NS or BH binaries is diabled here, since it is treated with KZ(27) now. \\
        & $\ge 1$: magnetic braking for NS \& BH (\texttt{brake.f}, \texttt{brake3.f})\\
        & $\ge 2$: Diagnostics at main output (\texttt{brake.f}) \\
        & $=3$: Input of ZMH = 1/SQRT(2*N) (Need KZ(5)$\ge$6) (\texttt{setup.F}) \\
        & $=4$: Set every star as type 13 Neutron star (Need KZ(27)=3) (\texttt{instar.f})  \\
KZ(29)  & (Reserved for new pulsar treatment decisions, not yet used)  \\
KZ(30)  & Hierarchical system regularization \\
        & $=-1$: Use chain only \\
        & $=0$: No triple, quad and chain regularization, only merger \\
        & $=1$: Use triple, quad and chain (\texttt{impact.f}) \\
        & $\ge 2$: Diagnostics at begin/end of chain at main output \\
        & $\ge 3$: Diagnostics at each step of chain at main output \\
KZ(31)  & Centre of mass correction after energy check (\texttt{cmcorr.f}) \\
KZ(32)  & Adjustment (increase) of adjust interval DTADJ, output interval DELTAT and energy error criterion QE based on binding eneryg of cluster (\texttt{check.f}) \\
KZ(33)  & Block-step statistics at main output (diagnostics) \\
        & $\ge 1$: Output irregular block step; and K.S. binary step if KZ(8)>0 \\
        & $\ge 2$: Output regular block step \\
KZ(34)  & Roche-lobe overflow \\
        & $=1$: Roche \& Spin synchronisation on binary with circular orbit (\texttt{synch.f}) \\
        & $=2$: Roche \& Tidal synchronisation on binary with circular orbit by BSE method (\texttt{bsetid.f}) \\
KZ(35)  & TIME reset to zero every 100 time units, total time is TTOT = TIME + TOFF (\texttt{offset.f})\\
KZ(36)  & (Suppressed) Step reduction for hierarchical systems \\
KZ(37)  & Neighbour list additions (\texttt{checkl.F}) \\
        & $\ge 1$: Add high-velocity particles into neighbor list \\
        & $\ge 2$: Add small time step particle (like close encounter particles near neighbor radius) into neighbor list \\
KZ(38)  & Force polynomial corrections during regular block step calculation \\
        & $=0$: no corrections \\
        & $=1$: all gains \& losses included \\
        & $=2$: small regular force change skipped \\
        & $=3$: fast neighbour loss only \\
KZ(39)  & Neighbor radius adjustment method \\
        & $=0$: The system has unique density centre and smooth density profile\\
        & $=1,\ge 3$: The system has no unique density centre or smooth density profile\\
        & ~~~~~~skip velocity modification of RS(I) (\texttt{regint.f}, \texttt{regcor\_gpu.f}) \\
        & ~~~~~~do not reduce neighbor radius if particle is outside half mass radius \\
        & ~~~~~~reduce RS(I) by multiply $0.9$ instead of estimation of RS(I) based on NNBOPT/NNB when neighbor list overflow happens (\texttt{fpoly0.F}, \texttt{util\_gpu.F}) \\
        & $=2,3$: Consider sqrt(particle mass / average mass) as the factor to determine the particle's neighbor membership. (\texttt{fpoly0.F}, \texttt{util\_gpu.F}) ({\bf Important to use if GPU is used for regular forces and neighbour lists!}) \\
KZ(40)  & $=0$: For the initialization of particle time steps, use only force and its first derivative to estimate. This is very efficent.\\
        & $>0$: Use fpoly2 (second and third order force derivatives calculation) to estimate the initial time steps. This method provides more accurate {\bf initial (only)} time steps and avoid incorrent time steps for some special cases like initially cold systems, but the computing cost (initial model only) is much higher ($O(N^2)$)\\
KZ(41)  & proto-star evolution of eccentricity and period for primordial binaries initialization (\texttt{proto\_star\_evol}, \texttt{binpop.F}) \\
KZ(42)  & Initial binary distribution \\
        & $=0$: RANGE>0: uniform distribution in log(semi) between SEMI0 and SEMI0/RANGE \\
        & ~~~~~~~ RANGE<0: uniform distribution in semi between SEMI0 and -1*RANGE. \\
        & $=1$: linearly increasing distribution function $f=0.03438*logP$ \\
        & $=2$: $f=3.5logP/[100+(logP)**2]$ \\
        & $=3$: $f=2.3(logP-1)/[45+(logP-1)**2]$; This is a ``3rd'' iteration when  pre-ms evolution is taken into account with KZ(41)=1 \\
        & $=4$: $f=2.5(logP-1)/[45+(logP-1)**2]$; This is a ``34th'' iteration when pre-ms evolution is taken into account with KZ(41)=1 and RBAR<$1.5$ \\
        & $=5$: Duquennoy \& Mayor 1991, Gaussian distribution with mean $\log{P}=4.8$, SDEV in $\log{P}=2.3$. Use Num.Recipes routine \texttt{gasdev.f} to obtain random deviates given ``idum1'' \\
%        & $=6$: eigen-evolution (Pavel Kroupa \& Rosemary Mardling)\\
KZ(43)  & Reserved for LPS+ code \url{https://github.com/kaiwu-astro/lps_plus} \\
KZ(44)  & (Unused) \\
KZ(45)  & $>0$: Computation of Moments of Inertia (with Chr. Theis) in output listing (\texttt{ellan.f}) \\
KZ(46)  &  HDF5/BINARY/ANSI format output and global parameter output (main output, see chapter~\ref{ch:output} for details) \\
        &  $=1,3$: HDF5(if HDF5 is compiled)/BINARY format\\
        &  $=2,4$: ANSI format \\
        &  $=1,2$: Only output active stars with time interval defined by KZ(47) \\
        &  $=3,4$: Output full particle list with time interval defined by KZ(47) \\
KZ(47)  &  Frequency for KZ(46) output. Should be non-negative integer. \\
        &  Output data with time interval $0.5^{KZ(47)} \times SMAX$ \\
KZ(48)  & $=0,1$: if 1 reduced HDF5 output (no 2nd and 3rd derivatives of force); warning: if KZ(40)=0 initial values wrong, undefined. \\
KZ(49)  & (Unused) \\
KZ(50)  & For event data bank, similar to MOCCA. Uses file units with numbers above 200; still under development; to switch on set to 1. \\

\end{longtable}

\hrule
\noindent
\begin{longtable}{@{}p{1.5cm}p{13.0cm}}
\caption{\texttt{Input of \&ININPUT, input.F line 8-9}}
\label{table:ininput89}\\\hline
%
DTMIN   & Time--step criterion for regularization search \\
RMIN    & Distance criterion for regularization search \\
ETAU    & Regularized time-step parameter (6.28/ETAU steps/orbit) \\
ECLOSE  & Binding energy per unit mass for hard binary (positive) \\
GMIN    & Relative two-body perturbation for unperturbed motion \\
GMAX    & Secondary termination parameter for soft KS binaries \\
SMAX    & Maximum time-step (factor of 2 commensurate with 1.0) \\
LEVEL   & \textcolor{red}{Stellar Evolution Level (see Kamlah et al. 2022), new input element} \\
\multicolumn{2}{l}{\textcolor{red}{\&INSSE, \&INBSE, \&INCOLL: new NAMELIST blocks read after LEVEL by subroutine READSE,}} \\
\multicolumn{2}{l}{\textcolor{red}{called from input.F; if blocks are empty in input, defaults according to LEVEL taken.}}
\end{longtable}

% & \texttt{if (kz(4).gt.0)} (suppressed now) \\
%DELTAS  & Output interval for binary search (in TCR; suppressed) \\
%ORBITS  & Minimum periods for binary output (level 1) \\
%GPRINT  & Perturbation thresholds for binary output (9 levels) \\

\hrule
\noindent
\begin{longtable}{@{}p{1.5cm}p{13.0cm}}
\caption{\texttt{Input of \&INDATA, data.F}}
\label{table:indata}\\\hline
%
ALPHAS  & Power-law index for initial mass function, routine \texttt{data.F}\\
BODY1   & Maximum particle mass before scaling (based on KZ(20); solar mass unit)\\
BODYN   & Minimum particle mass before scaling\\
NBIN0   & Number of primordial binaries (need KZ(8)>0) \\
        & -- by routine \texttt{imf2.F} using a binary IMF (KZ(20)$\ge$2)\\
        & -- by routine \texttt{binpop.F} splitting single stars (KZ(8)>0)\\
        & -- by reading subsystems from \texttt{dat.10} (KZ(22)$\ge$2)\\
NHI0    & Number of primordial hierarchical systems (need KZ(18)$\ge$2)\\
ZMET    & Metal abundance (in range 0.03 - 0.0001) \\
EPOCH0  & Formation time of stellar population, should be $\le$ 0, since at time of simulation start the stellar AGE is -EPOCH0 (in $10^6$ yrs); note that datsev.21 allows to input individual stellar ages \\
DTPLOT  & Plotting interval for stellar evolution HRDIAG (N-body units; $\ge$ DELTAT) \\
\end{longtable}

\hrule
\noindent
\begin{longtable}{@{}p{1.5cm}p{13.0cm}}
\caption{\texttt{Input of \&INSETUP, setup.F, KZ(5)>1}}
\label{table:insetup}\\\hline
       & \texttt{if (kz(5)=2)} \\\hline
APO    & Separation of two Plummer models in $N$--body units (SEMI = APO/(1 $+$ ECC). (Notice SEMI will be limited between $2.0$ and $50.0$) \\
ECC    & Eccentricity of two-body orbit (ECC $\ge$0 and ECC < $0.999$) \\
N2     & Membership of second Plummer model (N2 <= N) \\
SCALE  & Scale factor for the second Plummer model, second cluster will be generated by first Plummer model with $X \times SCALE$ and $V \times \sqrt{SCALE}$($\ge 0.2$ for limiting minimum size) \\\hline
       & \texttt{if (kz(5)=3)} \\\hline
APO    & Separation between the perturber and Sun in $N$--body units \\
ECC    & Eccentricity of orbit (=1 for parabolic encounter) \\
SCALE  & Perturber mass scale factor, perturber mass = Center star mass $\times$ SCALE (=1 for Msun) \\\hline
       & \texttt{if (kz(5)=4)} \\\hline
SEMI   & Semi-major axis (slightly modified; ignore if ECC > 1) \\
ECC    & Eccentricity (ECC > 1: NAME = 1 \& 2 free-floating) \\
M1     & Mass of first member (in units of mean mass) \\
M2     & Mass of second member (rescaled total mass = 1) \\\hline
       & \texttt{if (kz(5)$\ge$6) and (kz(24)<0)} \\\hline
ZMH    & Mass of single BH (in N-body units) \\
RCUT   & Radial cutoff in Zhao cusp distribution (MNRAS, 278, 488) \\
\end{longtable}

\hrule
\noindent
\begin{longtable}{@{}p{1.5cm}p{13.0cm}}
\caption{\texttt{Input of \&INSCALE, scale.F}}
\label{table:inscale}\\\hline
\texttt{scale.F:} & \\\hline
%
Q       & Virial ratio (routine \texttt{scale.F}; Q=0.5 for equilibrium) \\
VXROT   & XY--velocity scaling factor (> 0 for solid-body rotation) \\
VZROT   & Z--velocity scaling factor (not used if VXROT = 0) \\
RTIDE   & Unscaled tidal radius for KZ(14)=2 and KZ(22)$\ge$2. If not zero, RBAR = RT/RTIDE where RT[pc] is tidal radius calculated from input GMG and RG0\\
\end{longtable}

\hrule
\noindent
\begin{longtable}{@{}p{1.5cm}p{13.0cm}}
\caption{\texttt{Input of \&XTRNL0, xtrnl0.F}}
\label{table:indata}\\\hline
       & \texttt{if (kz(14)=2)} \\\hline
GMG    & Point-mass galaxy (solar masses, linearized tidel field in circular orbit) \\
RG0    & Central distance (in kpc) \\\hline
       & \texttt{if (kz(14)=3)} \\\hline
GMG    & Point-mass galaxy (solar masses) \\
DISK   & Mass of Miyamoto disk (solar masses) \\
A      & Softening length in Miyamoto potential (in kpc) \\
B      & Vertical softening length (kpc) \\
VCIRC  & Galactic circular velocity (km/sec) at RCIRC (=0: no halo) \\
RCIRC  & Central distance for VCIRC with logarithmic potential (kpc) \\
GMB    & Dehnen model budge mass (solar masses)\\
AR     & Dehnen model budge scaling radius (kpc)\\
GAM    & Dehnen model budge profile power index gamma \\  
RG(1:3) & Initial cluster position vector (kpc); DISK+VCIRC=0, VG(3)=0: A(1+E)=RG(1), E=RG(2) . Note that upon restart, RG is not read from input; saved value is used instead.\\
VG(1:3) & Initial cluster velocity vector (km/s). Note that upon restart, VG is not read from input; saved value is used instead. \\\hline
       & \texttt{if (kz(14)=3,4)} \\\hline
MP     & Total mass of Plummer sphere (in scaled units) \\
AP     & Plummer scale factor (N-body units; square saved in AP2) \\
MPDOT  & Decay time for gas expulsion (MP = MP0/(1 + MPDOT*(T-TD)) \\
TDELAY & Delay time for starting gas expulsion (T > TDELAY) \\\hline
       & \texttt{if (kz(14)=5)} \\\hline
RG(1:3) & Initial cluster position vector (kpc) . Note that upon restart, RG is not read from input; saved value is used instead.\\
VG(1:3) & Initial cluster velocity vector (km/s). Note that upon restart, VG is not read from input; saved value is used instead. \\ 
\end{longtable}

\hrule
\noindent
\begin{longtable}{@{}p{1.5cm}p{13.0cm}}
\caption{\texttt{Input of \&INBINPOP, binpop.F, \\\phantom{xxxxxxxx}if (kz(8)=1 or kz(8)>2)}}
\label{table:inbinpop}\\\hline
%
SEMI0   & Initial semi-major axis limit\\
ECC0    & Initial eccentricity \\
        & $<0$: thermal distribution, $f(e)=2e$ \\
        & $\ge 0$ and $\le 1$: fixed value of eccentricity\\
        & $=20$: uniform distribution \\
        & $=30$: distribution with $f(e)=0.1765/(e^2)$ \\
        & $=40$: general $f(e)=a*e^b$, $e0<=e<=1$ with $a=(1+b)/(1-e0^{(1+b)})$, current values: $e0=0$ and $b=1$ (thermal distribution) \\
RATIO   & KZ(42)$\le 1$: Binary mass ratio $M1/(M1 + M2)$\\
        & KZ(42)$=1.0$: $M1 = M2 = \langle M \rangle$\\
RANGE   & KZ(42)$=0$: semi-major axis range for uniform logarithmic distribution; \\
        & not used for other KZ(42)\\
NSKIP   & Binary frequency of mass spectrum (starting from body \#1)\\
IDORM   & Indicator for dormant binaries ($>0$: merged components)\\
\end{longtable}

\hrule
\noindent
\begin{longtable}{@{}p{1.5cm}p{13.0cm}}
\caption{\texttt{Input of \&INHIPOP, hipop.F\\\phantom{xxxxxxx} if (kz(8)>0 and kz(18)>1)}}
\label{table:inhipop}\\\hline
%
SEMI  &  Max semi-major axis in model units (all equal if RANGE = 0) \\
ECC   &  Initial eccentricity ($<0$ for thermal distribution) \\
RATIO &  Mass ratio ($= 1.0$: $M1 = M2$; random in [$0.5\thicksim0.9$]) \\
RANGE &  Range in SEMI for uniform logarithmic distribution ($> 0$) \\
\end{longtable}

\hrule
\noindent


%\hrule
%\noindent
%\begin{longtable}{@{}p{1.5cm}p{13.0cm}}
%\texttt{intide.F:}& \texttt{if (kz(27)>0)} (currently suppressed)
%RSTAR &  Size of typical star in A.U. \\
%IMS   &  # idealized main-sequence stars \\
%IEV   &  # idealized evolved stars \\
%RMS   &  Scale factor for main-sequence radii (>0: fudge factor) \\
%REV   &  Scale factor for evolved radii (initial size RSTAR)\\
%\end{longtable}

\hrule
\noindent
\begin{longtable}{@{}p{1.5cm}p{13.0cm}}
\texttt{cloud0.F:}& \texttt{if (kz(13)>0)} 
\textcolor{red}{Input currently NOT working with Namelist input method}
\\\hline
%
NCL   &  Number of interstellar clouds \\
RB2   &  Radius of cloud boundary in pc (square is saved) \\
VCL   &  Mean cloud velocity in km/sec \\
SIGMA &  Velocity dispersion (KZ(13)>1: Gaussian) \\
CLM   &  Individual cloud masses in solar masses (maximum MCL) \\
RCL2  &  Half-mass radii of clouds in pc (square is saved) \\
\end{longtable}

\vspace{1cm}

\section*{Stellar Evolution}

Stellar evolution is invoked by KZ(19)=1,2 or KZ(19)$\ge$3, offering two different
schemes.
The simpler one is KZ(19)=1, while the more complex one, K(19)$\ge$3,
is based on the Cambridge stellar evolution package
(Hurley, Pols, Tout 2000). The common envelope, roche transfering binaries are also considered. 
The main effects are changing stellar masses, radii, and luminosities,
which give rise to cluster mass loss.
The mass is assumed to escape from the cluster immediately and
possible collisions depend on stellar radii.

With the additional option KZ(12)>0, information on binaries and
single stars is written on two files (for binaries, unit 82, file \texttt{bev.82}; for single stars,
unit 83, file \texttt{sev.83}) in regular time intervals determined by
TPLOT (See details in Section~Output).
{\bf need to check for possibly more important information from Kamlah paper we missed out.}


\section*{Fortran NAMELIST based input}

The new input file is printed below. All input is grouped into a NAMELIST input block; our input is grouped according to the subroutines which read the lines ({\tt nbody6, input, data, setup, scale, xtrnl0, binpop, hipop}). In addition to that single and stellar evolution parameters, and stellar collision parameters, are read from the NAMELIST input blocks {\tt sse, bse, coll}. For every such NAMELIST input block the values can be given in the form {\tt name=value}, where {\tt name} is the variable name in the code. An input block in the file has to begin with {\tt \&NML} and terminate with {\tt /}, where {\tt NML} is the name of the block used in the code and in the input data file (e.g. {\tt INNBODY6}. Line breaks can be put in at liberty.

The principle of Fortran NAMELIST input is that the sequence in which the data are given is not relevant; if a variable of an input block is not given, it is possible to define it in the code (give default values). This principle is currently used for the input blocks {\tt sse, bse, coll} - a new variable {\tt LEVEL} of type {\tt CHARACTER*1} for the stellar evolution level is read at the end of the {\tt input} block - it can be Level A,B,C, or 0. Stellar Evolution Levels A,B,C are defined in Kamlah et al. (2022). If {\bf no data are given in} {\tt sse,bse,coll} {\bf the default values defined in Table A1 of that paper are used} (for Levels A,B,C, respectively). Level {\tt 0} (``zero'') corresponds to a run without stellar evolution. Note that {\tt LEVEL='0'} has to be provided in the input file for the case of no stellar evolution.

All variables can be redefined individually, for the first run and for any restart alike. For example
{\tt \&INSSE nsflag=3 /} in the input file will redefine nsflag different from the default. {\bf Notice: for all input blocks other than} {\tt sse, bse, coll} {\bf there are currently NO default values defined, so the user has to be responsible to provide a value for every required input} (as it was before). Values for all data in {\tt sse, bse, coll} are now always printed in the main output listing.


\begin{center}
\textcolor{red}{New Input File Example using Fortran NAMELIST}
\label{table:inputfile}\\
\fbox{\parbox{12.0cm}
{\tt
\&INNBODY6 \\
KSTART=1,TCOMP=1.0E8,TCRTP0=1000.0,\\isernb=40,iserreg=40,iserreg=0 / \\
\\
\&ININPUT \\
N=105000,NFIX=1,NCRIT=10,NRAND=43532,NNBOPT=80,NRUN=1,NCOMM=1, \\
ETAI=0.01,ETAR=0.01,RS0=0.15,DTADJ=1.0,DELTAT=1.0,\\TCRIT=1000.0,QE=1.0,RBAR=1.223,ZMBAR=0.7, \\
KZ(1:10)= 1 1 1 0 1 0 3 2 3 2 \\
KZ(11:20)=0 1 0 2 2 0 0 0 3 6 \\
KZ(21:30)=1 2 2 0 2 2 3 2 0 2 \\
KZ(31:40)=1 0 2 2 1 0 1 1 2 1 \\
KZ(41:50)=0 0 0 0 0 1 3 0 1 0 , \\
DTMIN=2.5E-6,RMIN=8.E-5,ETAU=0.1,ECLOSE=1.0,\\GMIN=1.0E-06,GMAX=0.01,SMAX=1.0, \\
Level='C' / \\
\\
\&INSSE / \\
\\
\&INBSE / \\
\\
\&INCOLL / \\
\\
\&INDATA \\
ALPHAS=2.35,BODY1=150.0,BODYN=0.08,NBIN0=5000,NHI0=0,\\ZMET=0.001,EPOCH0=0,DTPLOT=1.0 / \\
\\
\&INSETUP SEMI=,ECC=,APO=,N2=,SCALE=,ZM1=,ZM2,ZMH,RCUT= / \\
\\
\&INSCALE \\
Q=0.5,VXROT=0.0,VZROT=0.0,RTIDE=0.0 / \\
\\
\&INXTRNL0 \\
GMG=7.54E10,RG0=5.19,DISK=,A=,B=,VCIRC=,RCIRC=,GMB=,\\AR=,GAM=,RG=,,,VG=,,,MP=,AP2=,MPDOT=,TDELAY= / \\
\\
\&INHIPOP /
}
}
\end{center}


Note that empty placeholders like {\tt DISK=,} can be ignored, they just signal for the interested user where to put values if special options are chosen (as given in Table~\ref{table:allinput} and following tables).
\newpage

In the following all variables of the input blocks {\tt sse, bse, coll} are given, which are not displayed in the box above:

\begin{center}
\fbox{\parbox{12.0cm}
{\tt
\&INSSE \\
ecflag,wdflag,nsflag,psflag,mdflag,bhflag,kmech,mch,
mxns0,mxns1,nwind,bwind,flbv,disp,ecsig,wdsig1,wdsig2,wdkmax,\\
vvfac / \\
\&INBSE \\
lambd1,alphac,bhspinfl,kicktype,xk2,xk3,acc1,acc2,
xbeta,xxi,epsnov,eddfac,gamm1 / \\
\&INCOLL  \\
fctorcl /
}
}
\end{center}

A few notes: some variables had to be renamed compared to their genuine name in a certain subrouttine, in order to be unique throughout the code. These are:

\begin{center}
\fbox{\parbox{12.0cm}
{\tt
nwind = neta in mlwind.f \\
vvfac = vfac in kick.F \\
lambd1 = lambda in comennv.f \\
alphac = alpha1 in comenv.f \\
xk2 = k2 in hrdiag.f \\
xk3, xbeta, xxi = k3, beta, xi in mdot.F and roche.f 
}
}
\end{center}

The following variables are not explicitly documented in Kamlah et al. (2022), but included into the new namelist blocks: {\tt mch, mxns0, mxns}. These are the Chandrasekhar mass and two levels of neutron star masses; their defaults are defined equally for {\bf ALL} levels of stellar evolution as: {\tt mch=1.44, mxns0=1.8, mxns1 = 2.5} (solar masses).

Furthermore {\tt xk2, xk3, acc1} are not documented in the cited paper; their default values are set as {\tt xk2 = xk3 = 0.21, acc1 = 3.920659d8}; again these defaults do {\bf NOT} depend on stellar evolution and are taken from predefined values in the code before.
The value of {\tt k2 (xk2)} had been earlier hard-coded to be 0.21 in {\tt hrdiag.f}; now it can be set in input, the value from the input data is still used in {\tt hrdiag.f} to define {\tt k2}. {\tt hrdiag.f} argument list contains {\tt k2}, which is used in several other routines. 

The variables in the new input blocks {\tt sse, bse, coll} have been allocated now in three new Fortran {\tt common} blocks, with respective names. Due to that in several routines declarations of variables had to be adjusted, these are: \\[0.5cm]
\noindent
{\tt mlwind, mdot, comenv, corerd, roche, kick, kickGW, hrdiag, coal, mix}
 \\[0.5cm]
For handling the new input style using Fortran {\tt NAMELIST} input the following routines needed adjustment: \\[0.5cm]
\noindent
{\tt nbody6, input, modify, data, scale, xtrnl0, binpop, hipop}
\\[0.5cm]
The file {\tt input.F} contains a new subroutine {\tt READSE}, which handles read and write for all data in blocks {\tt sse, bse, coll}; both for initial run (called from {\tt input}) and for any restart (called from {\tt modify}.

\textcolor{red}{
Note that at the current time (April 2023) input done from the routinesIch 
{\tt cloud0, hotsys, ttinit, imbhinit }
{\bf fails! It is not working, and may be obsolete or updated in the future.}}


\section*{The old input file style (obsolete since April 2023)}

This input file was used until April 2023.
A typical input file can look like as follows.
It defines a new simulation running for 1,000,000 CPU--minutes
with $N=16,000$ particles distributed from a Plummer profile
(KZ(5)=1).
The run may alternatively terminate when TCRIT=1000.0 $N$--body units.
or if the final particle number of NCRIT=10 has been reached.
The output and adjustment time interval DELTAT/DTADJ are 1.0 N-body unit.
The initial mass function follows Kroupa, (2001) with mass ranging
from $m_{\rm max}=20.0 M_{\odot}$ to $m_{\rm min}=0.08 M_{\odot}$ (BODY1 and BODYN).
The initial virial ratio is 0.5 (equilibrium). The stellar 
evolution is switched on (KZ(19)=3) and initial metallicity is 0.001. 
Multiples and chain regularization are switched on (KZ(15)=2 and KZ(30)=2).
It uses solar neighbor tidal field (KZ(14)=1). 

\vspace{\baselineskip}
\noindent
\begin{center}
\fbox{\parbox{12.0cm}
{\tt
1 1000000.0 1.E6 40 40 640 \\
16000 1 10 43532 100 1 1 \\
0.02 0.02 0.1 1.0 1.0 1000.0 2.0E-05 1.0 0.7 \\
0 1 1 0 1 0 4 0 0 2 \\
0 1 0 1 2 1 0 0 3 6 \\
1 0 2 0 0 2 0 0 0 2 \\
1 0 2 1 1 0 1 1 0 0 \\
0 0 0 0 0 0 0 0 0 0 \\
1.0E-06 1E-4 0.2 1.0 1.0E-06 0.01 1.0 \\
2.35 20.0 0.08 0 0 0.001 0 1.0 \\
0.5 0.0 0.0 0.0 \\
7.54E10 5.19 \\
0.0005 -1.0 1.0 5.0 5 0 \\
\\
KSTART, TCOMP, TCRTP0, isernb,iserreg,iserks                  (nbody6.F)\\
N, NFIX, NCRIT, NRAND, NNBOPT, NRUN, NCOMM                    (input.F)
ETAI, ETAR, RS0, DTADJ, DELTAT, TCRIT, QE, RBAR, ZMBAR        (input.F)\\
(KZ(J),J=1,50)                                                (input.F)\\
DTMIN, RMIN, ETAU, ECLOSE, GMIN, GMAX, SMAX                   (input.F)\\
ALPHA, BODY1, BODYN, NBIN0, NHI0, ZMET, EPOCH0, DTPLOT        (data.F)\\
Q, VXROT, VZROT, RSPH2                                        (scale.F)\\
NBIN, SEMI0, ECC0, RATIO,RANGE, NSKIP, IDORM                  (binpop.f)\\
GMG, RG0                                                      (xtrnl0.F) KZ(14)=2 \\
GMG, DISK, A, B, VCIRC, RCIRC, GMB, AR, GAM,RG(1-3),VG(1-3)   (xtrnl0.F) KZ(14)=3 \\
SEMI0, ECC0, RATIO,RANGE, NSKIP, IDORM                        (binpop.f) KZ(8)=1,KZ(8)=3 \\
SEMI0, ECC0, RATIO,RANGE                                      (hipop.F) KZ(8)>0,KZ(18)>1
}
}
\end{center}


\section*{Input variables for primordial Binaries}

Many star clusters contain initial hard binaries with binding
energies much larger than the thermal energy
(the threshold ECLOSE is a suitable division between hard and
soft binaries).
There are two ways to initialise primordial binaries:

{\it The first} one always starts from some initial mass function (IMF)
provided by the routines \texttt{imf.f} or \texttt{imf2.f}.
The option KZ(8)=1 or $\ge 3$ invokes the routine \texttt{binpop.F}, which reads the
last line of the input file containing NBIN and the parameters of
their distribution (see above).
In this case, binaries are created either by random pairing of
single stars obtained from the IMF or by splitting them, depending on
the value of KZ(20) (see above).

{\it The second} way assumes that particle data, including the binaries,
are provided via the input data on file $\,$\texttt{dat.10}$\,$ (as e.g.
in the Kyoto--II collaborative experiment).
In such a case KZ(8)=2 and NBIN0 should be set to the expected number of primordial binaries
from the file. The code will first create NBIN0 centers of masses, and then use those
for scaling, before regularizing the pairs and the calculation begins.

A typical input file with primordial binaries looks as follows.
Here, we use binary random pairing from \texttt{imf2.f} and \texttt{binpop.F}
(KZ(20)=6 and KZ(8)=3, respectively) for 1000 initial binaries.
The semi-major axes of binaries use uniform distribution in log(semi) with a range from $41.3$ AU to $0.00413$ AU.
The eccentricity of binaries use thermal distribution. 
It was created from this input file running for 1000 time units.
Stellar evolution was also switched on in this file (KZ(19)=3).
In the package of the code, the file \texttt{N10k\_B1k.input} is
included.
%
%\vspace{\baselineskip}
%\noindent
\begin{center}
\fbox{\parbox{12.0cm}
{\tt
1 1000000.0 1.E6 40 40 640 \\
10000 1 10 43532 100 1 \\
0.02 0.02 0.17 1.0 1.0 800.0 5.0E-05 1.0 0.7 \\
0 1 1 0 1 0 4 3 0 2 \\
0 1 0 1 2 1 0 0 3 6 \\
1 0 2 0 0 2 0 0 0 2 \\
1 0 2 1 1 0 1 1 0 0 \\
0 0 0 0 0 0 0 0 0 0 \\
5.0E-06 3E-4 0.2 1.0 1.0E-06 0.01 0.5 \\
2.35 100.0 0.08 1000 0 0.001 0 1.0 \\
0.5 0.0 0.0 0.0 \\
2E-4 -1.0 1.0 1E4 5 0 
}
}
\end{center}



\vspace{0.7cm}

\section*{Restart}

\textcolor{red}{
Note that the new fortran NAMELIST based input does {\bf NOT} make any differences between input files for a restart or for a new run. The only difference will be the KSTART variable, which can bei either 1 or 2. All other values of KSTART will become obsolete; see next subsection.
}

It's very common that in the computer cluster every job has running time limit,
or the simulation stop due to some energy conservation problem or the normal stop
when the stop criterion is reached. In this case the user may want to continue the
simulation from the last time point. Thus the input parameter should changed to 
restart mode. The first line of input shown above combined together with two 
extra lines (See the description of KSTART in the parameter table above).
A simple example is :

\begin{center}
\fbox{\parbox{12.0cm}
{\tt
2 1000000.0 1.E6 40 40 640 
}
}
\end{center}

Here KSTART=2 means every parameter keeps the same value as before and just restarts from
the last saved file \texttt{comm.1}. If the user wants to change some parameters
of simulation, KSTART=3,5 can be set. For example:

\begin{center}
\fbox{\parbox{12.0cm}
{\tt
3 1000000.0 1.E6 40 40 640 \\
2.0 2.0 0.0 0.0 0.0 0.0 16 0 0.0 0.0
}
}
\end{center}

This restart file will change DTADJ and DELTAT to $2.0$. The KZ(16) is changed to 0.
All other parameters that are set to $0.0$ (TADJ, TNEXT, TCRIT, QE) keep same as before.
